\documentclass[11pt]{article}

\usepackage[margin=1in]{geometry}
\usepackage{booktabs}
\usepackage{graphicx}
\usepackage{amsmath}
\usepackage{siunitx}
\usepackage{hyperref}

\hypersetup{
  colorlinks=true,
  linkcolor=blue,
  citecolor=blue,
  urlcolor=blue
}

\title{Feature-First ISR: Semantic Transport for\\
Tactical UAS over Constrained and Contested Links}
\author{FFISR Project Team}
\date{\today}

\begin{document}
\maketitle

\begin{abstract}
Small UAS operating in contested environments routinely lose actionable
imagery at exactly the wrong moment: degraded links, RF blackouts, and
narrow spectrum windows compress or eliminate the video feed before it
reaches the operator. Feature-First ISR (FFISR) attacks this problem at
the transport layer. Rather than streaming continuous video, FFISR runs
object detection and multi-object tracking at the edge, transmits
compact semantic state at high cadence, and delivers annotated keyframe
imagery adaptively. Evaluated across 10 VisDrone sequences under three
link regimes, FFISR achieves \textbf{96\% bandwidth reduction} relative
to continuous video while maintaining semantic delivery under conditions
that produce 100\% baseline packet loss. The detector and tracker stack
is built on RT-DETR \cite{rtdetr2023} and ByteTrack or BoT-SORT
\cite{bytetrack2022,botsort2022}.
\end{abstract}

\section{Problem}
The ISR data chain for small UAS has a well-known weak point: the RF
link. Sensors have improved dramatically in resolution and frame rate,
edge compute is now powerful enough to run real-time neural detectors,
and ground stations can absorb high-rate data when conditions permit.
But the link itself is narrow, contested, and intermittent. When
throughput drops or signal is lost, a video-first pipeline delivers
nothing---the operator loses situational picture at the moment it
matters most.

Standard responses---compress harder, increase transmit power, add
spectrum---have real limits in tactical and regulatory contexts. A
different framing is: instead of fighting to keep the video moving,
prioritize the information the operator actually needs.

\section{Approach}

FFISR reorders the transport priority. Object-level state---track
identities, classes, positions, velocities, and confidence---is small,
structured, and high-value. A complete track update for a ten-object
scene fits in a few hundred bytes, surviving link conditions that would
corrupt or drop even a single compressed video frame. Full imagery is
transmitted only when it adds value: on operator request, on anomaly
detection, and on a link-adaptive periodic schedule.

Within the evaluation framework, a simulated baseline video channel
runs on the same emulated link under identical budget constraints.
This channel does not exist in the live deployment path; it is present
solely to give evaluators a direct comparison between traditional
streaming and the feature-first approach. The approach is consistent
with task-oriented communication research, which argues that
source-complete transmission is the wrong optimization target when the
receiver has a specific task \cite{semantic_survey_2021}.

\section{System Architecture}

\subsection{Edge Perception}

The perception layer runs on the same platform that carries the sensor.
A configurable detector---YOLO-family checkpoints or RT-DETR
\cite{rtdetr2023}---processes each frame and passes detections to a
tracker. ByteTrack \cite{bytetrack2022} and BoT-SORT
\cite{botsort2022} are both supported. The tracker maintains
persistent track identities across frames, producing per-update records
of track ID, class, confidence, bounding box, and estimated velocity.

RT-DETR is particularly well-suited here: unlike earlier
transformer-based detectors \cite{detr2020}, it achieves competitive
accuracy at real-time rates, viable on embedded UAS payloads.
ByteTrack's association strategy---keeping low-confidence detections in
play---improves continuity in cluttered, multi-scale aerial scenes
\cite{bytetrack2022}.

\subsection{Semantic Transport and Blackout Handling}

Each track update is serialized as compact JSON with delta-encoded
bounding boxes, keeping per-message size small for stable tracks and
allowing the channel to sustain high update cadence even under tight
bandwidth. The transport layer is link-agnostic: messages are written
to a configurable outbound queue and dispatched over whatever radio
interface the platform exposes, with no assumptions about underlying
protocol.

When tracker confidence drops below a threshold, the system emits an
anomaly event and optionally schedules a keyframe---giving the operator
an explicit signal when the perception layer is uncertain rather than
silently degrading.

RF blackouts are handled gracefully: outbound messages are buffered
locally and replayed on reconnect before resuming live updates. The
receiver reconstructs the current scene without restarting from zero,
and track IDs are preserved across the gap so objects reappear under
the same identities when link is restored.

\subsection{Adaptive Keyframe Policy}

Keyframes are full annotated frames transmitted over the semantic
channel on three triggers: explicit operator request, anomaly event,
and a periodic timer whose interval contracts under tighter link
budgets. Under 200~kbps the interval shortens to keep visual context
current. Keyframes are annotated with current track overlays so the
operator always receives context-aligned imagery.

\section{Metrics}

Throughput and savings figures are derived from delivery counters in
the transport layer. Let $B_s^{\text{att}}$ and $B_b^{\text{att}}$
denote bytes \emph{attempted} by the semantic and baseline channels
respectively, and $B_s$, $B_b$ the corresponding delivered bytes over
window $T$:

\begin{align}
\text{Sem.\ Throughput} &= \frac{8 B_s}{1000\,T} \quad (\text{kbps}) \\
\text{Base.\ Throughput} &= \frac{8 B_b}{1000\,T} \quad (\text{kbps}) \\
\text{BW Savings} &= 100\!\left(1 - \frac{B_s^{\text{att}}}{B_b^{\text{att}}}\right) \quad (\%)
\end{align}

Bandwidth savings are computed against attempted traffic to remain
valid even when the baseline channel is fully congested. A
transmission-volume proxy (total attempted megabits) captures the
relative data burden of each approach independent of link headroom.

\section{Evaluation}

Evaluation uses 10 sequences from the VisDrone benchmark
\cite{visdrone2019} with a fixed YOLOv8s detector and ByteTrack
tracker. Three link regimes cover the operational range from
permissive to severely constrained: Setting~A at 5000~kbps (no
blackout), Setting~B at 1000~kbps (5~s blackout window), and
Setting~C at 200~kbps (10~s blackout window). Packet loss is fixed
at 2\% across all settings. The detector and tracker configuration
is held constant---no per-sequence or per-regime tuning.

\subsection{Results}

\input{../../eval/results/summary_table.tex}

Table~\ref{tab:results} summarises the averaged results. Bandwidth
savings are consistent at $\sim$96\% across all regimes because the
semantic representation is compact regardless of link budget.
The resilience advantage is shown in Figure~\ref{fig:drops}: as the
link degrades from Setting A to C, baseline packet drop rises from
10\% to 100\%, while semantic drop remains below 4\%. At 200~kbps the
baseline channel delivers zero bytes; the semantic channel continues
operating normally.

Figure~\ref{fig:throughput} shows delivered versus required throughput
per regime. The semantic channel fits comfortably within the 200~kbps
budget ($\sim$100~kbps delivered); full-video streaming requires
2000--4600~kbps---an order of magnitude more.

\begin{figure}[ht]
\centering
\includegraphics[width=\columnwidth]{../../eval/results/fig_drops.pdf}
\caption{Packet drop rate as the link degrades. Baseline (grey) reaches
100\% at 200~kbps; semantic (blue) stays below 4\% across all
regimes. X-axis is inverted to show degrading conditions left-to-right.}
\label{fig:drops}
\end{figure}

\begin{figure}[ht]
\centering
\includegraphics[width=\columnwidth]{../../eval/results/fig_throughput.pdf}
\caption{Needed vs.\ delivered throughput per link regime. Hatched bars
show required bandwidth; solid bars show what was actually delivered.
Semantic messages fit within every regime; baseline video is shed
progressively as the budget tightens.}
\label{fig:throughput}
\end{figure}

\section{Conclusion}

FFISR demonstrates that a semantic-first transport policy maintains
tactical situational awareness under link conditions where continuous
video delivery fails completely. The software stack inserts at the
transport layer with no changes to sensor, radio, or ground station
hardware. The detector and tracker are configurable, so a program can
bring its own domain-trained model checkpoint. Acquisition stakeholders
retain a direct comparison channel throughout, which is a prerequisite
for building confidence in any AI-augmented system.

As spectrum competition increases and EW threats mature, the ability
to maintain tactical awareness on a degraded link becomes a force
multiplier. With $\sim$96\% bandwidth reduction and near-zero semantic
drop even at 200~kbps, FFISR represents a near-term, low-integration-risk
path to more resilient ISR on constrained tactical links. Detection
accuracy depends on training-domain match and may need adaptation for
IR or low-light operation; these are planned near-term extensions
alongside age-of-information metrics and expanded adverse-condition
trials.

\bibliographystyle{IEEEtran}
\bibliography{references}

\end{document}
